% Brief (tikz-diagram-craft), register row ch4-33.
% Reader question: given a flagged set that still contains a few true
%   positives, how many fewer opens does adaptive group-testing need than
%   reading every flag, and at what density does that stop paying?
% Claim: the guaranteed advantage adv(d)=1/(d(2*ceil(log2(1/d))+4)) crosses
%   1 at flagged-set density d=1/12: below that density zoom guarantees a
%   win, above it flat inspection can no longer be beaten on the worst case.
% Evidence: adv(1/12)=1.00 (crossing); adv(10/2500)=12.5, the F=2500,k=10
%   configuration whose realized mean was 162.6 opens (15.3x, better than
%   the 12.5x guarantee); adv(32/4096)=7.11 at the F=4096,k=32 tightness
%   config, where the exact tight formula lands on 511 queries, measured
%   511 exactly [internal, zoom_bound_check.py, run this pass].
% Must distinguish: the crossing point (adv=1, breakeven) from the two
%   worked configurations; that adv(d) is a genuine step function (the
%   ceiling), not a smooth curve, so no interpolation is drawn between
%   its jumps.
% Grammar chosen: threshold/regime plot, log-x density against guaranteed
%   advantage, a step curve (const plot) with a breakeven guide at adv=1.
% Grammar rejected: a bar chart of adv at the two named (F,k) configurations
%   would state the two numbers but hide the crossing itself -- the
%   chapter's point that d<=1/12 is a real, continuous threshold, not a
%   property of the two worked examples alone.
% Five-second test: "at the chapter's own F=2500,k=10 point, is the density
%   left or right of the d=1/12 crossing, and by how much?"
\begin{figure}[H]
\centering
\pdfigurehue{pdteal}%
\begin{tikzpicture}[pd figure]
\begin{axis}[pd axis,
  width=7.8cm,height=5.2cm,
  xmode=log,
  xmin=0.002,xmax=0.4,ymin=0,ymax=23,
  xtick={0.002,0.004,0.00781,0.0833,0.4},
  xticklabels={0.002,0.004,1/128,1/12,0.4},
  ytick={0,5,10,20},
  xlabel={flagged-set density $d=k/F$ (log scale)},
  ylabel={guaranteed advantage $\mathrm{adv}(d)$},
  clip=false,axis on top]
  \addplot[pd warn series,line width=.9pt,domain=0.002:0.4,samples=2] {1};
  \node[pd label,anchor=south east] at (axis cs:0.4,1) {breakeven: adv $=1$};
  \addplot[pd focus rule,line width=1.05pt,const plot,domain=0.002:0.4,samples=140,smooth=false]
    {1/(x*(2*ceil(ln(1/x)/ln(2))+4))};
  \draw[pd guide] (axis cs:0.0833,0) -- (axis cs:0.0833,1);
  \node[pd focus datum] at (axis cs:0.0833,1) {};
  \node[pd focus label,anchor=south west] at (axis cs:0.10,4.5) {$d{=}1/12$};
  \node[pd datum] at (axis cs:0.004,12.5) {};
  \node[pd label,anchor=south west,text width=2.8cm,align=left] at (axis cs:0.0043,14.0)
    {$F{=}2500,k{=}10$: $12.5\times$};
  \node[pd datum] at (axis cs:0.00781,7.11) {};
  \node[pd label,anchor=north west,text width=2.5cm,align=left] at (axis cs:0.011,6.0)
    {$F{=}4096,k{=}32$: $7.1\times$};
\end{axis}
\end{tikzpicture}
\caption{The zoom-advantage theorem's guaranteed worst-case advantage
against flagged-set density $d=k/F$ (log scale), a genuine step function
from the ceiling: $\mathrm{adv}(d)=1/(d(2\lceil\log_2(1/d)\rceil+4))$. The
guarantee crosses $1$ at $d=1/12$: at any sparser density, adaptive halving
is guaranteed to beat flat inspection, and at any denser one the worst case
can lose. At the chapter's own $F{=}2500,k{=}10$ point ($d{=}0.004$) the
guarantee is $12.5\times$, comfortably left of the crossing, and the
harness's realized mean of $162.6$ opens is a better-than-guaranteed
$15.3\times$; the $F{=}4096,k{=}32$ tightness configuration ($d{=}1/128$)
guarantees $7.1\times$ and lands exactly on its own tight closed form, $511$
queries measured against $511$ predicted. [internal,
\texttt{zoom\_bound\_check.py}]}
\label{fig:zoom-advantage}
\end{figure}
